Call-graph-guided retrieval is a natural fit for LLM-based fault localisation:
a structured graph constrains the search space and can either be navigated via
tool calls or serialised as a text slice in the prompt.
We evaluate whether two call-graph representations of \sqlite~3.30.1---a
compile-time graph extracted with \clang~18.1.3 and a syntactic graph extracted
with tree-sitter---improve LLM fault-localisation accuracy for six benchmark
questions on two large language models (\qwen and \nemotron).
We compare a tool-using agent that navigates the graph via function calling,
a text-slice condition that injects a BFS subgraph into the user message, and a
no-context baseline, together with four tool-ablation variants.
Graph context rarely improves exact accuracy under the 15-turn baseline: of the
43 wrong answers across all eight graph conditions, 15 (35\,\%) stem from function-pointer
callback gaps that make the target unreachable along static call edges, and 25 (58\,\%)
name a different graph node.
Text slices match the 15-turn tool condition in exact accuracy (and are within one question in
top-3 coverage) at about ten to twelve times lower token cost for \qwen and sixteen to
twenty-three times lower for \nemotron.
Neither graph representation (compile-time vs.\ syntactic) demonstrated superior accuracy.
In sensitivity probes (designed after the baseline results were known), a 30-turn budget with a revised prompt allowed agents to sometimes locate callback-separated targets, at steep token expense.
